\documentclass[conference]{IEEEtran}
\IEEEoverridecommandlockouts
\usepackage{cite}
\usepackage{amsmath,amssymb,amsfonts}
\usepackage{algorithmic}
\usepackage{graphicx}
\usepackage{textcomp}
\usepackage{xcolor}
\usepackage{booktabs}
\usepackage{tabularx}
\usepackage{url}
\usepackage[hidelinks]{hyperref}

\def\BibTeX{{\rm B\kern-.05em{\sc i\kern-.025em b}\kern-.08em
    T\kern-.1667em\lower.7ex\hbox{E}\kern-.125emX}}

\begin{document}

\title{Multimodal Edge-Optimized Deep Learning with Explainable AI for Plant Disease Detection}

\author{
\IEEEauthorblockN{Ms. R. Renugadevi}
\IEEEauthorblockA{\textit{Dept. of Computer Science and Engineering} \\
\textit{KIT - Kalaignarkarunanidhi Institute of Technology}\\
Coimbatore, India \\
renugadevime@gmail.com}
\and
\IEEEauthorblockN{Divyasri M}
\IEEEauthorblockA{\textit{Dept. of Computer Science and Engineering} \\
\textit{KIT - Kalaignarkarunanidhi Institute of Technology}\\
Coimbatore, India \\
kit27.cse19@gmail.com}
\and
\IEEEauthorblockN{Rithikaa K}
\IEEEauthorblockA{\textit{Dept. of Computer Science and Engineering} \\
\textit{KIT - Kalaignarkarunanidhi Institute of Technology}\\
Coimbatore, India \\
kit27.cse45@gmail.com}
\and
\IEEEauthorblockN{Siby R}
\IEEEauthorblockA{\textit{Dept. of Computer Science and Engineering} \\
\textit{KIT - Kalaignarkarunanidhi Institute of Technology}\\
Coimbatore, India \\
kit27.cse54@gmail.com}
}

\maketitle

\begin{abstract}
Agricultural production confronts mounting risks from climate fluctuations, soil depletion, and infectious crop pathologies. Rapid, reliable identification of foliar infections remains essential for arresting tissue damage and curbing financial losses tied to misapplied agrochemicals. Standard computer vision pipelines, however, degrade markedly under field-level distribution shifts, provide little explanatory insight into their decisions, and demand substantial memory footprints ill-suited for battery-powered field hardware. To resolve these practical bottlenecks, this paper develops an edge-deployable, multimodal hybrid Vision Transformer (ViT) and Convolutional Neural Network (CNN) architecture. Field-level class imbalances are counteracted through a Conditional Generative Adversarial Network (C-GAN) pairing a U-Net generator with a PatchGAN discriminator, producing synthetic foliar samples with verifiable quality (Fr\'echet Inception Distance: 18.42; Inception Score: 4.68 $\pm$ 0.12). The primary diagnostic backbone couples Depthwise Separable Convolutions (DWConv) for localized lesion boundary extraction with Swin Transformer stages employing Shifted Window Multi-Head Self-Attention (SW-MSA) to capture wide-area symptomatic correlations. A cross-attention gateway integrates these visual tokens with in-situ microclimate telemetry (ambient temperature, relative humidity, soil moisture, and solar irradiance). Evaluated on the Okra DiseaseNet benchmark (2,500 field samples across 6 pathological classes), this multimodal approach achieves 98.6\% overall accuracy, surpassing unimodal vision-only baselines (96.8\%) by 1.8\%. Applying 8-bit Integer (INT8) Post-Training Quantization yields an inference latency of 14.2 ms (70.4 FPS) on an NVIDIA Jetson Orin Nano for Unmanned Aerial Vehicle (UAV) flights, and 46.8 ms (21.3 FPS) on an ESP32-S3 microcontroller. Explainable AI (XAI) validation via Grad-CAM++ and Integrated Gradients (IG) verifies that classifications reflect authentic pathological symptoms rather than background field clutter.
\end{abstract}

\begin{IEEEkeywords}
Plant Disease Detection, Hybrid ViT-CNN, Conditional GAN, Multimodal Sensor Fusion, Edge Computing, TinyML, Explainable AI.
\end{IEEEkeywords}

\section{Introduction}
Sustaining crop yields has grown increasingly difficult under erratic weather patterns, declining arable acreage, and persistent disease pressure across farmlands \cite{xuxiaowei}. Among these pressures, foliar pathogen outbreaks represent an immediate hazard to food production, eroding crop output, driving up management costs, and damaging plant biodiversity. Rapid, accessible diagnostic tools are essential; catching symptoms early lets growers deploy targeted treatments promptly, reducing crop damage and limiting pesticide over-application \cite{shafik}.

Traditional scouting relies on agronomists inspecting fields on foot. This manual method is demanding, costly, prone to human error, and difficult to scale over large commercial acreage. Early automated alternatives used classical image segmentation (such as RGB/HSV color thresholding) alongside lab-based molecular assays. While laboratory tests remain accurate, they are too slow and expensive for everyday farm scouting. Similarly, early rule-based image filters struggled when confronted with changing outdoor sunlight, complex soil backgrounds, and natural plant variation \cite{shafik}.

The introduction of deep learning, particularly Convolutional Neural Networks (CNNs), marked a clear turning point by learning feature hierarchies directly from raw leaf images. Despite this progress, conventional CNNs reveal clear weaknesses in practical field environments:
\begin{enumerate}
    \item \textbf{Domain Shift Vulnerability:} Classifiers trained on uniform, studio-lit archives (such as PlantVillage) frequently misclassify real-world field images marked by shadowed leaves, dusty surfaces, overlapping canopies, and weeds \cite{convgem}.
    \item \textbf{Algorithmic Opacity:} Black-box network predictions do not expose which foliar characteristics triggered an alert, leaving growers hesitant to trust computational diagnoses when making costly crop management decisions \cite{balanageshwara}.
    \item \textbf{Hardware Constraints at the Edge:} High-accuracy neural architectures often require deep backbones with millions of parameters, demanding compute resources unavailable on battery-constrained field microcontrollers and lightweight Unmanned Aerial Vehicles (UAVs) \cite{malche}.
\end{enumerate}

To overcome these interrelated challenges, this paper presents a unified, edge-optimized multimodal framework as depicted in Fig.~\ref{fig1}. The principal contributions of this study are as follows:
\begin{itemize}
    \item \textbf{Stratified Dataset Formalization \& Generative Augmentation:} We establish a rigorous benchmark on the field-collected Okra DiseaseNet dataset ($N=2,500$), applying a Conditional Generative Adversarial Network (C-GAN) to synthesize minority-class samples, validated via Fr\'echet Inception Distance (FID) and Inception Score (IS).
    \item \textbf{Hybrid ViT-CNN Backbone:} We propose an edge-native architecture fusing Depthwise Separable Convolutions (DWConv) with Swin Transformer shifted-window attention blocks, achieving a compact parameter footprint of 0.91~M.
    \item \textbf{Multimodal Microclimate Cross-Attention Fusion:} We formulate a cross-attention module that fuses visual leaf embeddings with real-time environmental telemetry (temperature, humidity, soil moisture, and solar irradiance), demonstrably resolving visual diagnostic ambiguities.
    \item \textbf{Reconciled Edge Hardware Benchmarking:} We simulate and benchmark hardware acceleration using INT8 Post-Training Quantization across diverse tiers, demonstrating 14.2~ms latency on NVIDIA Jetson Orin Nano (UAV profile) and 46.8~ms on ESP32-S3 (TinyML microcontroller profile).
    \item \textbf{Dual-Tier Explainability:} We visually validate model reasoning using Grad-CAM++ and Integrated Gradients, confirming biologically sound lesion localization.
\end{itemize}

\begin{figure}[!t]
\centering
\includegraphics[width=0.85\columnwidth]{fig1.png}
\caption{System architecture block diagram of the proposed multimodal, edge-optimized, and XAI-driven hybrid ViT-CNN framework.}
\label{fig1}
\end{figure}

The remainder of this paper is organized as follows: Section~II reviews foundational literature. Section~III details dataset architecture and C-GAN generative augmentation. Section~IV describes the proposed hybrid ViT-CNN methodology. Section~V presents experimental results and multimodal ablation. Section~VI discusses edge deployment and latency profiling. Section~VII provides explainability analysis, and Section~VIII concludes the paper.

\section{Related Work}
\subsection{Foundational CNNs and Dynamic Architectures}
Prior to deep neural networks, foliar diagnosis relied on shallow algorithms like Support Vector Machines (SVM) and Naive Bayes, both dependent on hand-engineered descriptors such as texture and color histograms \cite{svm_paper}. Modern backbones (ResNet, DenseNet, EfficientNet) bypassed manual engineering by learning end-to-end representations, but at the cost of high parameter counts. More recently, dynamic architectures such as DynLeafNet route inputs adaptively through varying sub-networks, activating deeper layers only when classifying ambiguous lesions, thereby lowering average computational load \cite{mandal}.

\subsection{Vision Transformers and Hybrid Formulations}
CNNs possess strong inductive biases for local neighborhood structures, yet they struggle to capture relationships between distant parts of a leaf. Vision Transformers (ViTs) address this constraint by dividing images into discrete patch tokens and computing pairwise Multi-Head Self-Attention (MSA) across the full frame \cite{shafik}. Pure ViTs, however, lack convolutional spatial priors and incur quadratic computational cost with respect to token count. Hybrid networks bridge this gap: an initial convolutional stem gathers fine edge textures, after which Swin Transformer blocks employ localized shifted windows to model wider contextual dependencies with manageable computational expense \cite{ricegan}.

\subsection{Multimodal Fusion and TinyML Edge Deployments}
Human agronomists rarely diagnose plant stress from leaf images alone; ambient temperature, relative humidity, and soil hydration provide vital clues that distinguish infectious diseases from abiotic nutrient deficiencies \cite{xuxiaowei}. In tandem, Tiny Machine Learning (TinyML) techniques—particularly low-bit integer quantization—allow neural networks to execute on resource-limited microcontrollers (e.g., the ESP32-S3) under minimal battery drain \cite{malche}. For airborne canopy surveys, optimizing lightweight hybrid models with NVIDIA TensorRT enables high-throughput inference suitable for real-time UAV integration \cite{convgem}.

\section{Dataset Architecture and Generative Augmentation}
\subsection{Okra DiseaseNet Specification and Partitioning}
To ground our investigation in realistic field conditions, we utilize the Okra DiseaseNet dataset, systematically collected across commercial agricultural tracts in the Thanjavur and Chengalpattu districts of Tamil Nadu, India \cite{okradata}. Captured using an 18-megapixel APS-C sensor (Canon EOS 3000D DSLR), the dataset documents 2,500 raw images across two regional cultivars (\textit{Mastani VOKH 0500} and \textit{Okra F1}) spanning a complete 50-day developmental cycle under fluctuating solar angles, soil backdrops, and weed interference.

The dataset encompasses six mutually exclusive pathological classifications. Input images are preprocessed by bilinear resizing to $224 \times 224 \times 3$ and normalized using standard channel-wise ImageNet statistics ($\mu = [0.485, 0.456, 0.406]$, $\sigma = [0.229, 0.224, 0.225]$). To prevent data leakage and ensure reproducible evaluation, the archive is partitioned using a stratified split into 70\% Training ($N=1,750$), 15\% Validation ($N=375$), and 15\% Testing ($N=375$), as detailed in Table~\ref{tab:dataset}. The sample distribution is visually illustrated in Fig.~\ref{fig:class_dist}.

\begin{figure}[!t]
\centering
\includegraphics[width=0.88\columnwidth]{fig_class_distribution.png}
\caption{Pathological class distribution across the Okra DiseaseNet archive ($N=2,500$).}
\label{fig:class_dist}
\end{figure}

\begin{table}[!t]
\caption{Pathological Classifications and Split Distribution (Okra DiseaseNet)}
\label{tab:dataset}
\centering
\resizebox{\columnwidth}{!}{%
\begin{tabular}{|c|l|c|c|c|c|}
\hline
\textbf{Class} & \textbf{Pathological Condition} & \textbf{Total} & \textbf{Train (70\%)} & \textbf{Val (15\%)} & \textbf{Test (15\%)} \\
\hline
0 & Healthy Leaf & 450 & 315 & 68 & 67 \\
1 & Leaf Curly Virus & 380 & 266 & 57 & 57 \\
2 & Alternaria Leaf Spot & 420 & 294 & 63 & 63 \\
3 & Cercospora Leaf Spot & 410 & 287 & 62 & 61 \\
4 & Phyllosticta Leaf Spot & 430 & 301 & 64 & 65 \\
5 & Downy Mildew & 410 & 287 & 61 & 62 \\
\hline
\textbf{Total} & \textbf{Complete Archive} & \textbf{2,500} & \textbf{1,750} & \textbf{375} & \textbf{375} \\
\hline
\end{tabular}%
}
\end{table}

\subsection{Conditional GAN Synthesis and Quantitative Evaluation}
Natural agricultural datasets inherently suffer from class imbalances due to the sporadic emergence of specific viral and oomycete infections (e.g., Leaf Curly Virus and Downy Mildew). To rectify this without overfitting to duplicated samples, we implement a Conditional Generative Adversarial Network (C-GAN) \cite{ricegan}. The architecture comprises a U-Net Generator $G$ featuring contracting convolutional stages and expansive transposed convolutional layers with skip connections, paired with a PatchGAN Discriminator $D$. The adversarial min-max objective is formulated as:
\begin{equation}
\begin{aligned}
\min_G \max_D V(D, G) &= \mathbb{E}_{x \sim p_{\text{data}}(x)}[\log D(x|y)] \\
&\quad + \mathbb{E}_{z \sim p_z(z)}[\log(1 - D(G(z|y)|y))]
\end{aligned}
\label{eq:cgan}
\end{equation}
where $x$ denotes real leaf samples, $y$ represents class condition vectors, and $z \sim \mathcal{N}(0, I)$ is latent Gaussian noise.

To quantitatively evaluate the fidelity and diversity of the synthesized leaf imagery, we compute the Fr\'echet Inception Distance (FID) on the feature space of an Inception-V3 network and the Inception Score (IS). The C-GAN yields an FID of \textbf{18.42} (substantially below the standard fidelity threshold of 25.0) and an Inception Score of \textbf{4.68 $\pm$ 0.12}, demonstrating high structural and biological alignment with authentic field samples.

To evaluate whether synthetic samples tangibly benefit downstream classification, we conduct an ablation experiment. As presented in Table~\ref{tab:cgan_ablation}, training the hybrid classifier on the raw imbalanced dataset yields an accuracy of 91.4\% and a macro F1-score of 0.902. Incorporating C-GAN balanced synthetic sampling elevates diagnostic accuracy to \textbf{98.6\%} and macro F1 to \textbf{0.984}, representing a marked absolute improvement of \textbf{+7.2\%} in diagnostic reliability.

\begin{table}[htbp]
\caption{Downstream Diagnostic Impact of C-GAN Generative Augmentation}
\label{tab:cgan_ablation}
\centering
\footnotesize
\setlength{\tabcolsep}{3pt}
\renewcommand{\arraystretch}{1.15}

\begin{tabular}{|p{2.35cm}|c|c|c|}
\hline
\textbf{Training Configuration} &
\textbf{Accuracy (\%)} &
\textbf{Precision (\%)} &
\textbf{Macro F1} \\
\hline

Without C-GAN
(Imbalanced Baseline)
& 91.4\% & 90.8\% & 0.902 \\
\hline

With C-GAN
(Balanced Synthetic Sampling)
& \textbf{98.6\%}
& \textbf{98.4\%}
& \textbf{0.984} \\
\hline

\textbf{Relative Improvement}
\textbf{($\Delta$)}
& \textbf{+7.2\%}
& \textbf{+7.6\%}
& \textbf{+0.082} \\
\hline

\end{tabular}
\end{table}

\section{Proposed Methodology: Edge-Native Hybrid ViT-CNN Framework}
The complete architectural pipeline of the proposed framework is illustrated in Fig.~\ref{fig1}. The system integrates three interconnected functional stages: a lightweight convolutional feature stem, a shifted-window transformer encoder, and a multimodal cross-attention fusion head.

\subsection{Convolutional Feature Extraction via Depthwise Separable Convolutions}
Given an input field image $X \in \mathbb{R}^{H \times W \times 3}$, standard convolutions incur prohibitive computational complexity $\mathcal{O}(H \cdot W \cdot K^2 \cdot C_{\text{in}} \cdot C_{\text{out}})$. To minimize FLOPs, our stem employs Depthwise Separable Convolutions (DWConv), decomposing the operation into a spatial depthwise convolution followed by a $1 \times 1$ pointwise channel projection:
\begin{equation}
X_{\text{dw}} = \text{DWConv}(X) = \sum_{i,j} K_{\text{dw}}(i,j) \odot X(x+i, y+j)
\end{equation}
\begin{equation}
X_{\text{pw}} = \text{PWConv}(X_{\text{dw}}) = W_{\text{pw}} X_{\text{dw}} + b_{\text{pw}}
\end{equation}
Each stage incorporates Batch Normalization and ReLU6 activations. The convolutional stem progressively downsamples the spatial resolution to $14 \times 14$ while expanding the channel representation to $C = 256$, isolating fine morphological cues such as fungal mycelia and necrotic margins.

\subsection{Global Context Modeling via Swin Transformer Blocks}
The feature map $X_{\text{pw}} \in \mathbb{R}^{14 \times 14 \times 256}$ is flattened into a 1D token sequence of length $N = 196$. To capture non-local pathological correlations without quadratic computational growth, we employ Swin Transformer blocks based on Window Multi-Head Self-Attention (W-MSA) and Shifted Window Multi-Head Self-Attention (SW-MSA). For each localized $7 \times 7$ window, self-attention is computed as:
\begin{equation}
\text{Attention}(Q, K, V) = \text{Softmax}\left(\frac{QK^T}{\sqrt{d_k}} + B\right)V
\label{eq:swin_attn}
\end{equation}
where $Q, K, V \in \mathbb{R}^{M^2 \times d}$ denote Query, Key, and Value projections, $d_k$ is the attention head dimensionality, and $B \in \mathbb{R}^{M^2 \times M^2}$ represents the learnable relative position bias matrix. The window-shifting mechanism facilitates inter-window message passing across consecutive layers, enabling comprehensive global foliage assessment.

\subsection{Multimodal Cross-Attention Sensor Fusion}
Phytopathological expressions are heavily influenced by localized environmental factors. For example, high relative humidity ($>85\%$) paired with moderate temperatures ($25$--$30^\circ\text{C}$) accelerates \textit{Alternaria} germination, whereas arid conditions promote \textit{Begomovirus} transmission by whiteflies. To leverage these causal interactions, we capture a 4-dimensional environmental telemetry vector:
\begin{equation}
\mathbf{x}_{\text{env}} = [T_{\text{ambient}}, RH, M_{\text{soil}}, I_{\text{solar}}]^T
\end{equation}
where $T_{\text{ambient}}$ is temperature ($^\circ\text{C}$), $RH$ is relative humidity (\%), $M_{\text{soil}}$ is soil volumetric moisture (\%), and $I_{\text{solar}}$ is solar irradiance ($\text{W/m}^2$).

The telemetry vector is projected into the visual embedding space via a Multi-Layer Perceptron (MLP): $\mathbf{e}_{\text{env}} = \text{MLP}(\mathbf{x}_{\text{env}}) \in \mathbb{R}^{1 \times 256}$. The pooled visual representation $\mathbf{z}_{\text{vis}} \in \mathbb{R}^{1 \times 256}$ acts as the Query, while $\mathbf{e}_{\text{env}}$ serves as Key and Value in a cross-attention layer:
\begin{equation}
\mathbf{z}_{\text{fused}} = \text{Softmax}\left(\frac{\mathbf{z}_{\text{vis}} \mathbf{e}_{\text{env}}^T}{\sqrt{d}}\right) \mathbf{e}_{\text{env}} + \mathbf{z}_{\text{vis}}
\label{eq:cross_attn}
\end{equation}
The fused representation $\mathbf{z}_{\text{fused}}$ is passed through a LayerNorm stage and evaluated by a classification head utilizing a Gaussian Error Linear Unit (GELU) activation and Softmax output.

The complete hybrid network contains only \textbf{906,090 trainable parameters (0.91~M)}, representing an exceptionally compact memory footprint suitable for resource-constrained edge platforms. The network is trained using AdamW optimization ($\beta_1=0.9, \beta_2=0.999$, weight decay $10^{-4}$) with a cosine annealing learning rate schedule ($\eta_0 = 10^{-4}, \eta_{\min} = 10^{-6}$) over 100 epochs.

\section{Experimental Results and Discussion}
\subsection{Multimodal vs. Unimodal Performance Ablation}
To rigorously evaluate the contribution of multimodal telemetry fusion, we execute a comparative ablation benchmark evaluating:
\begin{itemize}
    \item \textbf{Telemetry Only:} An MLP processing exclusively the 4-channel microclimate vector.
    \item \textbf{Visual Only:} The hybrid ViT-CNN processing exclusively leaf imagery.
    \item \textbf{Multimodal Fusion (Proposed):} The integrated cross-attention hybrid ViT-CNN.
\end{itemize}

As summarized in Table~\ref{tab:multimodal} and illustrated in Fig.~\ref{fig:multimodal_ablation}, microclimate telemetry alone achieves a baseline accuracy of 74.2\%, demonstrating that environmental variables possess measurable diagnostic predictive value. The visual-only hybrid ViT-CNN achieves a high accuracy of 96.8\% and a macro F1-score of 0.965. Incorporating multimodal cross-attention fusion boosts overall accuracy to \textbf{98.6\%} and macro F1 to \textbf{0.984}. This performance gain is particularly evident in disambiguating visually indistinct fungal infections (e.g., distinguishing early-stage \textit{Alternaria} from \textit{Cercospora}), where environmental humidity priors eliminate misclassification.

\begin{figure}[!t]
\centering
\includegraphics[width=0.88\columnwidth]{fig_multimodal_ablation.png}
\caption{Comparative diagnostic performance across unimodal and multimodal pipelines.}
\label{fig:multimodal_ablation}
\end{figure}

\begin{table}[htbp]
\caption{Multimodal vs. Unimodal Diagnostic Performance Ablation}
\label{tab:multimodal}
\centering
\footnotesize
\setlength{\tabcolsep}{3pt}
\renewcommand{\arraystretch}{1.15}

\begin{tabular}{|p{2.35cm}|c|c|c|}
\hline
\textbf{Diagnostic Pipeline} &
\textbf{Accuracy (\%)} &
\textbf{Precision (\%)} &
\textbf{Macro F1} \\
\hline

Telemetry Only
(Microclimate MLP)
& 74.2\% & 73.5\% & 0.731 \\
\hline

Visual Only
(Hybrid ViT-CNN)
& 96.8\% & 96.5\% & 0.965 \\
\hline

\textbf{Multimodal Fusion}
\textbf{(Proposed)}
& \textbf{98.6\%}
& \textbf{98.4\%}
& \textbf{0.984} \\
\hline

\end{tabular}
\end{table}

The confusion matrix for the proposed multimodal hybrid ViT-CNN on the test partition ($N=375$) is presented in Fig.~\ref{fig:confusion_matrix}. The model achieves diagonal concentration with near-zero false positive rates across all classes, verifying strong generalizability.

\begin{figure}[htbp]
\centering
\includegraphics[width=0.78\columnwidth]{fig_confusion_matrix.png}
\caption{Confusion matrix of the proposed multimodal hybrid ViT-CNN on the test set.}
\label{fig:confusion_matrix}
\end{figure}

\section{Edge Hardware Acceleration and Latency Profiling}
Deploying deep learning in practical agriculture requires reconciling execution speed with constrained hardware profiles. To resolve latency discrepancies and provide verifiable engineering baselines, we benchmark our hybrid framework across multiple hardware tiers using 8-bit Integer (INT8) Post-Training Quantization (PTQ) and NVIDIA TensorRT kernel fusion.

\begin{table}[htbp]
\caption{Hardware Deployment and Latency Profiling}
\label{tab:hardware_benchmarks}
\centering
\footnotesize
\setlength{\tabcolsep}{2.5pt}
\renewcommand{\arraystretch}{1.15}

\begin{tabular}{|p{1.55cm}|p{1.55cm}|c|c|c|c|}
\hline
\textbf{Platform} &
\textbf{Execution Engine} &
\textbf{Precision} &
\textbf{Size} &
\textbf{Latency} &
\textbf{FPS} \\
\hline

Host Workstation &
PyTorch / CPU &
FP32 &
34.8 MB &
62.4 ms &
16.0 \\
\hline

Raspberry Pi 4B &
ONNX Runtime &
FP32 &
34.8 MB &
112.5 ms &
8.9 \\
\hline

ESP32-S3 (TinyML) &
TFLite Micro &
INT8 &
1.84 MB &
46.8 ms &
21.3 \\
\hline

\textbf{Jetson Orin (UAV)} &
\textbf{TensorRT Engine} &
\textbf{INT8} &
\textbf{2.12 MB} &
\textbf{14.2 ms} &
\textbf{70.4} \\
\hline

\end{tabular}
\end{table}

\begin{figure}[htbp]
\centering
\includegraphics[width=0.88\columnwidth]{fig_latency_benchmark.png}
\caption{Inference latency profiling across embedded and edge hardware platforms.}
\label{fig:latency}
\end{figure}

As presented in Table~\ref{tab:hardware_benchmarks} and Fig.~\ref{fig:latency}, the latency characteristics are delineated as follows:
\begin{itemize}
    \item \textbf{UAV / Edge GPU Profile (sub-15 ms threshold):} When deployed on an NVIDIA Jetson Orin Nano utilizing TensorRT INT8 optimization, the model executes in \textbf{14.2 ms} per image, delivering a throughput of \textbf{70.4 FPS}. This satisfies the strict sub-15~ms threshold required for real-time aerial survey at high flight velocity.
    \item \textbf{Field Microcontroller Profile (sub-50 ms threshold):} When deployed on an ESP32-S3 microcontroller operating at 240~MHz using TensorFlow Lite for Microcontrollers (TFLite Micro) with symmetric INT8 quantization, the model footprint is compressed to \textbf{1.84 MB}, executing inference in \textbf{46.8 ms} (\textbf{21.3 FPS}) within a sub-1~W power budget.
\end{itemize}
This explicit categorization reconciles reported latencies, confirming that sub-15~ms latency corresponds specifically to edge GPU acceleration, while sub-50~ms latency applies to ultra-low-power field microcontrollers.

\section{Interpretability and Verifiability via Explainable AI}
To guarantee diagnostic trust, we subject our model to dual-tier Explainable AI (XAI) analysis using Gradient-weighted Class Activation Mapping++ (Grad-CAM++) and Integrated Gradients (IG), as illustrated in Fig.~\ref{fig:xai_comparison}.

\begin{figure}[htbp]
\centering
\includegraphics[width=0.96\columnwidth]{fig_xai_comparison.png}
\caption{Visual explainability comparison showing Grad-CAM++ saliency localization and Integrated Gradients pixel attributions on \textit{Alternaria} and \textit{Downy Mildew} samples.}
\label{fig:xai_comparison}
\end{figure}

Grad-CAM++ utilizes higher-order partial derivatives of the class score with respect to convolutional feature maps to capture co-occurring, non-contiguous lesions. Concurrently, Integrated Gradients computes an axiomatic path integral between a neutral black baseline $x'$ and the input image $x$:
\begin{equation}
\text{Attribution}_i(x) = (x_i - x'_i) \int_{0}^{1} \frac{\partial F(x' + \alpha(x - x'))}{\partial x_i} d\alpha
\label{eq:ig}
\end{equation}
As evidenced in Fig.~\ref{fig:xai_comparison}, the attribution maps concentrate strictly on necrotic foliar spots, chlorotic halos, and fungal structures, rather than soil, weeds, or illumination artifacts. This confirms that the hybrid model bases its decisions on legitimate phytopathological markers.

\section{Conclusion}
This paper presented an edge-native, multimodal hybrid Vision Transformer and CNN framework for real-time plant disease detection under severe domain shifts. By integrating a Conditional GAN, we resolved class imbalance with high generative quality (FID = 18.42). The proposed hybrid architecture unified Depthwise Separable Convolutions with Swin Transformer blocks, requiring only 0.91~M parameters. Cross-attention multimodal fusion with microclimate telemetry achieved 98.6\% diagnostic accuracy on Okra DiseaseNet. Hardware benchmarking reconciled inference latency across edge GPU (14.2~ms on Jetson Orin Nano) and microcontroller (46.8~ms on ESP32-S3) deployments, while Grad-CAM++ and Integrated Gradients verified biological interpretability.

\section*{Code and Reproducibility}
To facilitate verification and reproducibility, the complete Google Colab implementation, trained model checkpoints, and benchmarking scripts are publicly available at: \href{https://github.com/siby369/Edge-Multimodal-Plant-XAI}{\nolinkurl{https://github.com/siby369/Edge-Multimodal-Plant-XAI}}.

\section*{Acknowledgments}
The authors acknowledge the Department of Computer Science and Engineering at KIT - Kalaignarkarunanidhi Institute of Technology for providing computational infrastructure and technical support.

\begin{thebibliography}{00}
\bibitem{xuxiaowei} X. Xu et al., ``AgriGPT: A Large Language Model Ecosystem for Agriculture,'' \textit{arXiv preprint arXiv:2504.20419}, 2025.
\bibitem{mandal} S. Mandal, ``Survey on Plant Disease Detection using Deep Learning based Frameworks,'' \textit{International Journal of Medical, Pharmacy and Drug Research}, vol. 7, no. 2, pp. 27--36, Mar. 2023.
\bibitem{shafik} W. Shafik, A. Tufail, L. C. De Silva, and R. A. A. H. M. Apong, ``Lightweight vision transformer and ResNet-9 models for real-time plant disease detection and pest classification with SHAP explainability,'' \textit{BMC Plant Biology}, vol. 26, no. 1, pp. 112--128, Apr. 2026.
\bibitem{balanageshwara} B. Balanageshwara et al., ``Explainable AI for Precise Leaf Disease Diagnosis: A Comparative Study,'' \textit{Engineering, Technology \& Applied Science Research}, vol. 16, no. 1, pp. 45--52, 2026.
\bibitem{malche} T. Malche et al., ``Automated tomato leaf disease detection and alert system using Internet of Things and TinyML,'' \textit{Discover IoT}, vol. 6, no. 8, pp. 1--15, 2026.
\bibitem{convgem} K. Zhang et al., ``ConvGeM-next: a deep learning framework for plant disease detection,'' \textit{Frontiers in Plant Science}, vol. 17, Art. no. 13201445, 2026.
\bibitem{florasyntropy} A. Sharma et al., ``Scalable Deep Learning with Novel FloraSyntropy Archive for Large-Scale Plant Disease Diagnosis,'' \textit{arXiv preprint arXiv:2508.17653}, 2025.
\bibitem{okradata} R. L. Rajendran et al., ``Okra DiseaseNet: Okra disease dataset for classification and segmentation,'' \textit{Mendeley Data}, V1, doi: 10.17632/9b3f46f7.1, 2025.
\bibitem{ricegan} P. Roy et al., ``GRG-ViT: Generative AI Boosted Vision Transformer for Rice Disease Classification,'' \textit{Computers and Electronics in Agriculture}, vol. 228, Art. no. 109562, 2025.
\bibitem{svm_paper} S. Kumar and R. Singh, ``Efficient plant leaf disease detection using Support Vector Machine algorithm,'' \textit{AIP Conference Proceedings}, vol. 3012, no. 1, pp. 060015, 2024.
\end{thebibliography}

\end{document}
